\documentclass[twocolumn]{article}
\usepackage[utf8]{inputenc}
\usepackage{amsmath, amssymb, amsfonts}
\usepackage{graphicx}
\usepackage{booktabs}
\usepackage{hyperref}
\usepackage{algorithm}
\usepackage{algpseudocode}

\title{Unsupervised Domain Adaptation for Traditional Mongolian Historical Archives Recognition via Adversarial Stroke Enhancement}
\author{Your Name, AI Assistant}
\date{\today}

\begin{document}

\maketitle

\begin{abstract}
The digitization of traditional Mongolian historical archives is severely hindered by the scarcity of labeled data and the extreme degradation of paper materials. In this paper, we propose an unsupervised domain adaptation (UDA) framework specifically designed for archival manuscripts. Our approach combines a style-aligned synthetic engine with a Domain Adversarial Neural Network (DANN). To address the fading ink problem in historical archives, we introduce a Multi-Scale Stroke Enhancement (MSSE) module that adaptively recalibrates stroke features. Experimental results on real archival documents demonstrate that our method effectively bridges the gap between synthetic fonts and historical handwriting without requiring any human annotations.
\end{abstract}

\section{Introduction}
Traditional Mongolian archives represent a vital part of World Cultural Heritage. However, the automated recognition of these documents (HTR) remains an open challenge due to: (1) The complex vertical orientation and cursive ligatures. (2) The severe "Archival Noise" including ink bleeding, paper yellowing, and mechanical damage. (3) The total absence of large-scale ground truth labels for historical hands.

\section{Proposed Method}
\subsection{Synthesis for Archival Realism}
We generate a source domain $D_s$ by rendering Unicode text on real archival textures extracted from original TIFF files. An elastic transformation $\mathcal{G}$ is applied to simulate the variance of archival handwriting.

\subsection{Adversarial Domain Adaptation}
To align the synthetic source $D_s$ and the unlabeled archival target $D_t$, we employ a minimax game between a feature extractor $G_f$ and a domain discriminator $G_d$. The Gradient Reversal Layer (GRL) ensures that the learned features are invariant to the specific archival background noise.

\subsection{Multi-Scale Stroke Enhancement (MSSE)}
Given the fading nature of archival ink, MSSE utilizes dilated convolutions to extract multi-scale contextual information of strokes, ensuring that even faint connections in cursive script are captured.

\section{Experimental Evaluation}
\subsection{Datasets}
\begin{itemize}
    \item \textbf{Source}: 2,320 synthetic Mongolian columns with archival style.
    \item \textbf{Target}: Real unlabeled columns from historical archives (19th-century administrative records).
\end{itemize}

\subsection{Results Analysis}
The unsupervised model achieved significant qualitative improvements. As shown in our pseudo-label results, the model successfully identifies keywords such as "Khuree" (ᠬᠣᠠᠠᠨ) which were previously misrecognized by baseline models.

\section{Conclusion}
This work provides the first end-to-end unsupervised pipeline for traditional Mongolian archives recognition. By leveraging adversarial learning and archival-style synthesis, we achieve a breakthrough in recognizing documents for which manual labeling is economically or technically infeasible.

\end{document}
